\documentclass[12pt]{article}

% ------------------------------------------------------------
% Packages
% ------------------------------------------------------------
\usepackage{amsmath, amssymb, amsthm, bm}
\usepackage{geometry}
\usepackage{setspace}
\usepackage{hyperref}
\usepackage{physics}
\usepackage{mathtools}

\geometry{margin=1in}
\setstretch{1.15}

% ------------------------------------------------------------
% Theorem Environments
% ------------------------------------------------------------
\newtheorem{theorem}{Theorem}
\newtheorem{lemma}{Lemma}
\newtheorem{definition}{Definition}
\newtheorem{proposition}{Proposition}

% ------------------------------------------------------------
% Title
% ------------------------------------------------------------

\title{\textbf{Deep Learning Solvers as Lamphrodynamic Flows\\
in RSVP Field Theory}}

\author{Flyxion}

\date{}

% ------------------------------------------------------------
% Document
% ------------------------------------------------------------

\begin{document}

\maketitle

\begin{abstract}
Deep learning methods for solving high-dimensional dynamic economic 
models have recently gained attention for their ability to ``unbend'' 
complex solution manifolds and overcome the curse of dimensionality. 
This essay presents a unified geometric and variational interpretation 
of Jesús Fernández-Villaverde’s deep learning solvers through the lens 
of the Relativistic Scalar Vector Plenum (RSVP) field theory. By 
identifying value functions, policy rules, and informational structures 
with scalar potentials, vector flows, and entropy fields in RSVP, we 
show that deep networks implement a lamphrodynamic smoothing flow that 
discovers the natural coordinate charts of the economic equilibrium 
manifold. We derive a free-energy functional whose minimization is 
equivalent to satisfying equilibrium conditions, formalize the 
associated gradient flow, and prove equivalence theorems linking 
economic equilibrium with RSVP stationary field configurations.
\end{abstract}

% ============================================================
\section{Introduction}

In recent work, Fernández-Villaverde has argued that deep neural 
networks provide a powerful and flexible method for solving dynamic 
economic models, especially in high-dimensional settings where 
traditional numerical approaches become intractable. He emphasizes the 
geometric role of deep networks as machines that ``unbend'' complex 
high-dimensional solution manifolds into simpler latent representations.

This geometric perspective aligns naturally with the 
Relativistic Scalar Vector Plenum (RSVP), a field-theoretic framework 
in which economic objects are embedded as scalar, vector, and entropy 
fields living on a curved manifold. RSVP provides an intrinsic geometric 
and variational interpretation of deep learning solvers as 
\emph{lamphrodynamic flows}---curvature-smoothing, entropy-reducing 
transformations that evolve field configurations toward equilibrium.

This essay develops the full mapping between the economic, deep 
learning, and RSVP perspectives, culminating in a free-energy functional 
whose minimization yields an economic equilibrium.

% ============================================================
\section{Equilibrium as a Field Configuration}

Dynamic economic models require finding a function
\begin{equation}
    D(\mathbf{x})
\end{equation}
that satisfies an equilibrium operator equation:
\begin{equation}
    T(D) = \mathbf{0},
\end{equation}
where $D$ may represent a value function, policy rule, or both.

\subsection*{RSVP Field Mapping}

In RSVP, the unknown function $D$ corresponds to a tuple of fields:
\[
(\Phi, \mathbf{v}, S),
\]
where
\begin{itemize}
    \item $\Phi(\mathbf{x})$: a scalar potential (value function),
    \item $\mathbf{v}(\mathbf{x})$: a vector flow (policy dynamics),
    \item $S(\mathbf{x})$: an entropy field (uncertainty, beliefs).
\end{itemize}

Thus the equilibrium becomes a geometric constraint:
\begin{equation}
    T(\Phi, \mathbf{v}, S) = \mathbf{0}.
\end{equation}

\begin{definition}[RSVP Equilibrium Field]
A triple $(\Phi, \mathbf{v}, S)$ is an RSVP equilibrium if it satisfies
\[
T(\Phi,\mathbf{v},S)=0.
\]
\end{definition}

\noindent
This is analogous to a fixed point, zero-curvature state, or a stationary point of a variational energy.

% ============================================================
\section{Dimensionality as Curvature and Entropy}

The curse of dimensionality arises when the solution surface is highly 
curved or kinked. RSVP identifies:
\begin{itemize}
    \item \textbf{Curvature}: large $\|\nabla \Phi\|$ or $\|\nabla \mathbf{v}\|$,
    \item \textbf{Torsion}: nondifferentiabilities induced by constraints,
    \item \textbf{Entropy-of-representation}: irregular information 
    distribution in $S$.
\end{itemize}

Standard bases (e.g., polynomials) assume low curvature and low entropy. 
Deep networks, by contrast, learn transformations that reduce both.

% ============================================================
\section{Neural Networks as Lamphrodynamic Operators}

A single neural layer performs the map
\begin{equation}
    f_\theta(\mathbf{x}) = \sigma(W\mathbf{x}+b),
\end{equation}
a local nonlinear transformation smoothing curvature and concentrating 
entropy. RSVP identifies such a transformation as a 
\emph{lamphrodyne update}.

A deep network composed of $L$ such layers:
\[
\mathbf{x}
\stackrel{f_1}{\mapsto} \mathbf{x}_1
\stackrel{f_2}{\mapsto} \cdots
\stackrel{f_L}{\mapsto} \mathbf{x}_L
\]
approximates a continuous lamphrodynamic flow:
\begin{equation}
    \frac{d\mathbf{x}}{dt} = \mathcal{V}(\mathbf{x}),
\end{equation}
where $\mathcal{V}$ is a curvature-smoothing vector field.

% ============================================================
\section{A Free-Energy Functional for Economic Equilibrium}

To interpret training as variational descent, we introduce a 
free-energy functional.

\begin{definition}[RSVP Free-Energy Functional]
Let $D = D(\Phi,\mathbf{v},S)$. Define:
\begin{equation}
\mathcal{E}[D]
    = \frac{1}{2}\int_{\mathcal{X}}
        \|T(D(\mathbf{x}))\|^2 \, d\mu(\mathbf{x}).
\end{equation}
\end{definition}

We now prove that minimizing this energy is equivalent to satisfying the 
equilibrium operator.

\begin{lemma}[First Variation of $\mathcal{E}$]
Let $D_\epsilon = D + \epsilon \delta D$. Then
\[
\frac{d}{d\epsilon}\mathcal{E}[D_\epsilon]\bigg|_{\epsilon=0}
    = \int_{\mathcal{X}}
        T(D)\,
        \frac{\delta T}{\delta D}[\delta D]
        \, d\mu.
\]
\end{lemma}

\begin{proof}
Expand:
\[
\mathcal{E}[D_\epsilon]
= \frac{1}{2}\int \|T(D + \epsilon \delta D)\|^2\,d\mu.
\]
Differentiate under the integral:
\[
\frac{d}{d\epsilon}
\mathcal{E}[D_\epsilon]
= \int T(D) \cdot \frac{d}{d\epsilon}T(D+\epsilon\delta D)\, d\mu.
\]
At $\epsilon=0$, this yields the stated expression.
\end{proof}

\begin{theorem}[Equivalence of Minimizers and Equilibria]
A field configuration $D^\ast$ minimizes $\mathcal{E}$ if and only if
\[
T(D^\ast)=0.
\]
\end{theorem}

\begin{proof}
If $T(D^\ast)=0$, then $\mathcal{E}[D^\ast]=0$ and is clearly minimal.

Conversely, suppose $D^\ast$ minimizes $\mathcal{E}$. Then the first 
variation must vanish for all perturbations:
\[
\int T(D^\ast)\,
\frac{\delta T}{\delta D}[\delta D]\, d\mu = 0
\quad \forall\,\delta D.
\]
If $\delta T/\delta D$ is nondegenerate, this implies $T(D^\ast)=0$.
\end{proof}

\begin{proposition}
Gradient descent on $\mathcal{E}$ generates a lamphrodynamic flow
\begin{equation}
\begin{aligned}
\frac{d\Phi}{dt} &= - \frac{\delta \mathcal{E}}{\delta \Phi},\\
\frac{d\mathbf{v}}{dt} &= - \frac{\delta \mathcal{E}}{\delta \mathbf{v}},\\
\frac{dS}{dt} &= - \frac{\delta \mathcal{E}}{\delta S}.
\end{aligned}
\end{equation}
\end{proposition}

\begin{proof}
Immediate from functional gradient descent applied to 
$\mathcal{E}[\Phi, \mathbf{v}, S]$.
\end{proof}

% ============================================================
\section{Formal RSVP Formulation of Dynamic Economic Models}

Let $\mathcal{X}$ be the state manifold. Define:
\[
\Phi: \mathcal{X}\to\mathbb{R}, \qquad
\mathbf{v}: \mathcal{X}\to T\mathcal{X}, \qquad
S: \mathcal{X}\to\mathbb{R}_{\ge 0}.
\]

\begin{definition}[RSVP Economic Model]
A dynamic economic model is the tuple 
$(\mathcal{X},\Phi,\mathbf{v},S,T)$
satisfying
\[
T(\Phi,\mathbf{v},S)=0.
\]
\end{definition}

\begin{theorem}[RSVP Characterization of Equilibrium]
The economic equilibrium is the stationary point of the lamphrodynamic 
gradient flow on the free-energy functional:
\[
(\Phi^\ast,\mathbf{v}^\ast,S^\ast)
=
\arg\min_{(\Phi,\mathbf{v},S)}
    \mathcal{E}[\Phi,\mathbf{v},S].
\]
\end{theorem}

\begin{proof}
Direct from the previous theorem and definitions.
\end{proof}

% ============================================================
\section{Variational and Field-Theoretic Perspectives}

The lamphrodynamic interpretation of deep learning as a curvature-reducing flow
naturally invites a broader variational and field-theoretic perspective.
Within this conceptual frame, an economic equilibrium is not merely the fixed
point of an operator, but the stationary configuration of a collection of fields
whose interactions encode incentives, information, constraints, and adaptive
dynamics.

In the standard formulation, one seeks a function \(D(\mathbf{x})\)---a value
function, policy function, or other decision rule---that satisfies an operator
equation \(T(D)=0\).  In the geometric formulation adopted here, the same
object is represented by a triplet of fields
\[
(\Phi, \mathbf{v}, S),
\]
where \(\Phi\) denotes a scalar potential associated with accumulated value or
fitness, \(\mathbf{v}\) denotes a vector flow describing the directional
structure of optimal adjustment, and \(S\) denotes an entropy or informational
field induced by beliefs, dispersion, or uncertainty.

From this vantage, the equilibrium condition \(T(D)=0\) admits a natural
translation into a stationarity requirement for a free-energy-like functional
\(\mathcal{E}[\Phi,\mathbf{v},S]\).  The discrete steps of a numerical method,
including those implemented by deep learning procedures, can then be viewed as
approximations to a continuous relaxation dynamics in which the fields evolve
toward the stationary configuration:
\[
\frac{\delta \mathcal{E}}{\delta \Phi}
= \frac{\delta \mathcal{E}}{\delta \mathbf{v}}
= \frac{\delta \mathcal{E}}{\delta S}
= 0.
\]
This variational interpretation does not replace the computational perspective;
instead, it clarifies why certain numerical schemes succeed, why they produce
stable solutions, and why they behave as if they were performing implicit
geometric simplifications of the underlying manifold.  The economic equilibrium
emerges, in this light, as the minimal-action configuration of an interacting
field system.

% ============================================================
\section{Dynamics, Canonical Structure, and Information Flow}

Once the fields \((\Phi,\mathbf{v},S)\) are treated as dynamical quantities, it
is natural to inquire whether the relaxation toward equilibrium admits a
canonical structure.  The temporal evolution of these fields defines an infinite
dimensional phase space, in which each field possesses a conjugate momentum
derived from its role in the Lagrangian or free-energy functional.

In this representation, iterative learning procedures, policy iterations,
stochastic updates, and gradient-based optimization can all be interpreted as
discrete approximations to continuous flows generated by a Hamiltonian or
generalized Hamilton--Jacobi structure.  The vector field governing these flows
is not arbitrary: it encodes the propagation of information, the dissipation of
irregularity, and the alignment of beliefs and actions, in exactly the way that
canonical flows transmit forces, curvature, or probability mass in classical
and statistical field theories.

This canonical viewpoint distinguishes genuine economic adjustment directions
from redundant or gauge-like transformations of the fields.  In particular, it
reveals the extent to which variations in \(\mathbf{v}\) reflect actual
behavioral or decision-theoretic changes, while variations in \(S\) reflect
informational or representational pressures.  The interplay between these
components can then be treated as a structured flow that respects both the
geometry of the state manifold and the integrability constraints implied by
economic rationality.

This dynamical framework provides a unifying structure in which learning,
adaptation, inference, and optimization appear as different manifestations of a
single, coherent geometric process.

% ============================================================
\section{Scale Interactions and Global Inference}

Economic systems exhibit interactions across multiple scales: individual
decisions respond to local conditions and expectations, while aggregate
outcomes depend on global distributions and long-range informational
structure.  The field-theoretic representation makes these scale interactions
explicit by encoding them in the triplet \((\Phi,\mathbf{v},S)\).

The scalar potential \(\Phi\) captures smooth, low-frequency structure such as
persistent value gradients or intertemporal incentives.  The vector field
\(\mathbf{v}\) describes local adjustment dynamics and short-scale sensitivity.
The entropy field \(S\) records global dispersion, uncertainty, heterogeneity,
and belief formation.  Together, these fields determine how local optimization
criteria propagate, how global consistency is enforced, and how equilibrium
configurations aggregate information.

This structure admits a natural interpretation in terms of global inference
problems.  High-dimensional decision environments often contain latent
low-dimensional order, but this structure may not be visible in the raw state
space.  Deep learning solvers implicitly identify this order by performing
multi-scale smoothing and alignment steps, which correspond to reducing the
curvature and entropy of the fields while preserving the essential causal
relationships.

From this perspective, the success of deep learning methods in solving
high-dimensional equilibrium models is not accidental.  These methods perform a
sequence of transformations that resemble renormalization or coarse-graining:
they reduce representational complexity while retaining the structures relevant
for equilibrium.  The final solution is therefore best understood as the fixed
point of a multi-scale information flow, in which local incentives and global
constraints coexist coherently.

This multi-scale geometric perspective offers new tools for analyzing stability,
bifurcations, aggregation, and robustness in large economic systems, while
preserving the interpretability and normative structure of equilibrium
analysis.

\section{Equilibrium, Information, and Futarchic Governance}

If the lamphrodynamic interpretation of deep learning reveals the intrinsic
geometry of equilibrium, it also clarifies the informational structure that any
governance system must confront.  In particular, the RSVP formulation offers a
natural bridge to Hanson’s futarchy paradigm, in which policy selection is
guided by predictive markets whose prices encode beliefs about welfare-relevant
futures.

In this setting, the fields \((\Phi,\mathbf{v},S)\) provide a precise analogue
of the informational, behavioral, and uncertainty components that futarchy
seeks to aggregate.  The scalar potential \(\Phi\) represents an intertemporal
welfare or value landscape; the vector field \(\mathbf{v}\) captures feasible
adjustments and behavioral responses; and the entropy field \(S\) encodes the
dispersion of beliefs, expectations, and informational frictions.  Policy
evaluation becomes a question of how proposed interventions deform these
fields, alter the curvature of the welfare landscape, or redistribute entropy
across states.

From this perspective, prediction markets do not merely forecast scalar
aggregates; they express constraints on the allowable field configurations.
Market-implied expectations correspond to boundary data for the fields, and
policy comparison becomes the evaluation of alternative stationary solutions
under different boundary or coupling structures.  The futarchic principle
``vote on values, bet on beliefs'' thereby acquires a geometric interpretation:
collective preferences determine the potential to be optimized, while market
prices determine the informational terms that regulate how beliefs shape the
flow of the fields.

This integration suggests that futarchy is not an institutional curiosity but a
natural extension of a field-theoretic understanding of equilibrium.  The same
variational logic that governs economic dynamics, representation learning, and
lamphrodynamic smoothing also governs the aggregation of dispersed information
about the future.  A futarchic mechanism can therefore be seen as selecting,
from among many possible stationary field configurations, the one that
maximizes collective welfare subject to informational coherence.  This provides
a unified mathematical and conceptual foundation linking equilibrium
computation, predictive governance, and high-dimensional learning dynamics.

\section{A Formal Variational Theorem for Futarchic Policy Selection}

The geometric interpretation of futarchy can be made precise by formulating
policy choice as a constrained variational problem on the RSVP fields
\((\Phi,\mathbf{v},S)\).  Let \(\theta\) denote a feasible policy parameter, and
let the resulting economic environment induce a free-energy functional
\(\mathcal{E}_\theta[\Phi,\mathbf{v},S]\) whose stationary points correspond to
equilibrium field configurations.  Prediction markets generate
belief-weighted boundary data that constrain the admissible fields.  Let
\(\mathcal{B}_\theta\) denote this set of boundary conditions for policy
\(\theta\).

We now formalize the relationship between futarchy and equilibrium selection.

\begin{theorem}[Futarchic Stationarity Principle]
Suppose that for each feasible policy parameter \(\theta\), the free-energy
functional \(\mathcal{E}_\theta\) is convex in a neighborhood of its stationary
points and that the associated boundary set \(\mathcal{B}_\theta\) is
nonempty, closed, and defined by market-implied expectations of
welfare-relevant outcomes.  Then the futarchic mechanism selects the policy
\(\theta^\ast\) whose equilibrium field configuration
\((\Phi_\theta,\mathbf{v}_\theta,S_\theta)\) satisfies
\[
(\Phi_\theta,\mathbf{v}_\theta,S_\theta)
\in \argmin_{(\Phi,\mathbf{v},S)\in \mathcal{B}_\theta}
\mathcal{E}_\theta[\Phi,\mathbf{v},S],
\]
and
\[
\theta^\ast
\in \argmax_{\theta}
W\!\left[\Phi_\theta,\mathbf{v}_\theta,S_\theta\right],
\]
where \(W\) is a welfare functional encoding the value-judgment component of
collective preference.  Moreover, \(\theta^\ast\) is stable under perturbations
of the boundary data corresponding to small revisions in market beliefs.
\end{theorem}

\begin{proof}
For any fixed \(\theta\), convexity of \(\mathcal{E}_\theta\) in a neighborhood
of its stationary points ensures the existence of a unique minimizer subject to
\(\mathcal{B}_\theta\).  Since \(\mathcal{B}_\theta\) consists of boundary data
derived from prediction-market prices, the minimizer incorporates all
belief-weighted expectations about future states.

Let \((\Phi_\theta,\mathbf{v}_\theta,S_\theta)\) denote this constrained
minimizer.  By definition, it satisfies the variational stationarity conditions
\[
\frac{\delta \mathcal{E}_\theta}{\delta \Phi} = 0,
\qquad
\frac{\delta \mathcal{E}_\theta}{\delta \mathbf{v}} = 0,
\qquad
\frac{\delta \mathcal{E}_\theta}{\delta S} = 0
\qquad
\text{on the interior,}
\]
together with the boundary constraints specified by \(\mathcal{B}_\theta\).

Collective value judgments determine the welfare functional \(W\).  A
futarchic mechanism chooses the policy that maximizes expected welfare
evaluated at the equilibrium configuration.  Thus,
\[
\theta^\ast
= \argmax_\theta W\!\left[\Phi_\theta,\mathbf{v}_\theta,S_\theta\right].
\]

Finally, stability follows from the closedness of \(\mathcal{B}_\theta\) and
continuity of both the minimizer and the welfare functional in the underlying
boundary data.  Small perturbations in prediction-market prices induce small
changes in the boundary set, which in turn induce small variations in the
stationary solution and therefore do not alter the maximizing policy except
under degeneracy.
\end{proof}

This result shows that futarchy can be interpreted as a two-level variational
problem: fields evolve internally to satisfy equilibrium constraints, while
policy parameters are selected externally to maximize welfare under the
belief-weighted informational structure encoded by markets.  From this
perspective, futarchy implements a principled mechanism for selecting among
equilibrium field configurations under uncertainty, fully consistent with the
variational and geometric structures underlying the RSVP representation.


% ============================================================
\section*{Conclusion}

Deep learning solvers for economic models naturally implement a form of
lamphrodynamic flow in the sense of RSVP field theory.  Neural networks act as
automatic diffeomorphism machines: they bend, straighten, and reparameterize the
state space through successive nonlinear transformations that reduce curvature,
contract representational entropy, and reveal the intrinsic geometry of the
equilibrium manifold.  What appears in standard computational terms as feature
extraction or representation learning is, in geometric terms, an iterative
search for a coordinate system in which the economic fields
\((\Phi,\mathbf{v},S)\) admit a simple and nearly linear description.  The
network discovers, rather than assumes, the structure that renders the solution
tractable.

Training corresponds to a variational descent on a free-energy functional whose
minimizers coincide with the equilibria of the economic model.  Stochastic
gradient steps, random sampling, and iterative updates become discrete
approximations to a continuous relaxation process in which the fields evolve
toward the stationary configuration that satisfies all incentive, feasibility,
and informational constraints.  The monotonic improvement properties observed in
practice reflect the geometry of this descent: the flow aligns the scalar
potential with intertemporal rewards, the vector field with feasible and optimal
adjustment dynamics, and the entropy field with coherent aggregation of
uncertainty and beliefs.

This perspective coheres naturally with the broader variational and dynamic
framework developed in the preceding sections.  Once economic decision rules are
understood as field configurations, and once their numerical approximation is
interpreted as a curvature-smoothing process, previously disparate phenomena
acquire a unified explanation.  Stability of iterative solvers appears as
convexity or geodesic regularity of the free-energy landscape.  Generalization
emerges as the selection of low-entropy, globally coherent field
configurations.  Canonical structure becomes visible in the evolution of the
fields and their conjugate momenta.  Multi-scale effects—from individual
decisions to aggregate distributions—find representation in the interplay
between local curvature and global entropy.

Seen in this light, the integration of futarchic governance becomes natural
rather than extraneous.  Predictive markets supply belief-weighted boundary
conditions for the fields, while collective value judgments define the welfare
functional to be optimized.  The futarchic stationarity principle makes this
precise: among the family of stationary equilibria permitted by the boundary
data, the mechanism selects the policy whose induced field configuration
maximizes welfare subject to informational coherence.  Policy choice is thus
elevated from a comparative statics exercise to a variational selection problem
embedded directly within the geometry of the RSVP manifold.

The synthesis accomplished here therefore provides not merely an interpretation
of Fernández-Villaverde’s geometric intuition but a rigorous and extensible
field-theoretic framework unifying machine learning, equilibrium computation,
predictive governance, and high-dimensional economic dynamics.  The
lamphrodynamic viewpoint clarifies why deep learning succeeds where traditional
projection and perturbation methods fail, and it offers new tools for
understanding equilibrium selection, stability, bifurcation, robustness, and
policy design in complex economic environments.  At the same time, it reveals
deep structural affinities between economic reasoning, variational physics,
predictive markets, and modern representation learning: all seek, in their own
languages, the coordinate systems in which the dynamics become simple, the
fields become smooth, and the equilibrium becomes manifest.


\appendix
% ============================================================
% Appendix: Field-Theoretic Lagrangian for RSVP-Economics
% with Futarchy Coupling (Hanson-Inspired)
% ============================================================

\section*{Appendix: A Field-Theoretic Lagrangian for RSVP Economics\\
with Futarchy Coupling Mechanisms}

This appendix presents a formal Lagrangian formulation of dynamic
economic models within the Relativistic Scalar Vector Plenum (RSVP)
framework. We integrate Robin Hanson’s principle of \emph{futarchy}
---“vote on values, bet on beliefs”—as a set of gauge-type couplings
between value potentials, predictive-belief fields, and policy flows.
The result is a unified variational principle governing economic
equilibrium, informational aggregation, and decision selection.

% ------------------------------------------------------------
\subsection*{A. Field Content}

Let $\mathcal{X}$ be the economic state manifold. We define the RSVP
fields:
\[
\Phi(\mathbf{x},t): \text{scalar potential (value field)},\\[4pt]
\mathbf{v}(\mathbf{x},t): \text{policy vector field},\\[4pt]
S(\mathbf{x},t): \text{entropy / belief-distribution field},\\[4pt]
B(\mathbf{x},t): \text{futarchy prediction field (market-belief gauge)}.
\]

The field $B$ represents the market’s aggregated probabilistic beliefs
about future outcomes. Hanson’s futarchy asserts that $B$ should control
\emph{which policy} $\mathbf{v}$ is selected, while $\Phi$ encodes the
collective welfare target over which society “votes.”

% ------------------------------------------------------------
\subsection*{B. Futarchy Principles as Field Constraints}

Hanson’s motto—``vote on values, bet on beliefs’’—is expressed through:
\begin{enumerate}
    \item \textbf{Value Principle:}
        \[
        \Phi \text{ defines the welfare objective (value votes).}
        \]
    \item \textbf{Belief Principle:}
        \[
        B \text{ encodes predictive accuracy (market bets).}
        \]
    \item \textbf{Decision Principle:}
        \[
        \mathbf{v} \text{ is selected to maximize predicted } \Phi
        \text{ under } B.
        \]
\end{enumerate}

Mathematically, this induces a coupling between $B$ and $\mathbf{v}$
analogous to a gauge interaction.

% ------------------------------------------------------------
\subsection*{C. Free-Energy and Action Principles}

Let the RSVP free-energy functional be:
\[
\mathcal{F}[\Phi,\mathbf{v},S]
    = \frac{1}{2}\int_{\mathcal{X}}
        \|T(\Phi,\mathbf{v},S)\|^2 \, d\mu,
\]
where $T(\Phi,\mathbf{v},S)=0$ is the equilibrium operator.

We extend this by adding:
\begin{itemize}
\item a \textbf{dynamical term} for predictively accurate beliefs,
\item a \textbf{coupling} between beliefs $B$ and policy flows $\mathbf{v}$,
\item a \textbf{value alignment potential} coupling $\Phi$ to $B$.
\end{itemize}

The full \emph{RSVP-Futarchy Action} is:
\[
\mathcal{S} = \int dt \int_{\mathcal{X}} 
\mathcal{L}(\Phi, \mathbf{v}, S, B) \, d\mu.
\]

% ------------------------------------------------------------
\subsection*{D. The Lagrangian Density}

We propose the following Lagrangian:
\begin{equation}
\boxed{
\mathcal{L}
=
\mathcal{L}_{\text{dyn}}
+
\mathcal{L}_{\text{equil}}
+
\mathcal{L}_{\text{belief}}
+
\mathcal{L}_{\text{futarchy}}
+
\mathcal{L}_{\text{entropy}}
}
\end{equation}

Each term is defined as follows.

% -------------------
\subsubsection*{1. Dynamical Term}

Lamphrodynamic evolution of fields:
\[
\mathcal{L}_{\text{dyn}}
=
\frac{1}{2}\Big(
    |\partial_t \Phi|^2
  + \|\partial_t \mathbf{v}\|^2
  + |\partial_t S|^2
  + |\partial_t B|^2
\Big).
\]

% -------------------
\subsubsection*{2. Equilibrium Constraint Term}

Enforces the RSVP economic equilibrium via penalty (free-energy):
\[
\mathcal{L}_{\text{equil}}
=
- \frac{\lambda}{2}\,
    \|T(\Phi, \mathbf{v}, S)\|^2.
\]

Large $\lambda$ recovers exact equilibrium.

% -------------------
\subsubsection*{3. Belief Dynamics Term}

Prediction market accuracy emerges from minimizing surprisal:
\[
\mathcal{L}_{\text{belief}}
=
- \frac{\alpha}{2}
    \|\nabla B\|^2
    - \beta\, B\, S.
\]

Interpretation:
- $\|\nabla B\|^2$: smoother belief fields generalize better  
- $B S$: beliefs must compete with entropy/noise

% -------------------
\subsubsection*{4. Futarchy Coupling Term}

This encodes actions selected by beliefs to maximize values:
\[
\mathcal{L}_{\text{futarchy}}
=
\gamma\, B\,
    \langle \nabla \Phi,\ \mathbf{v} \rangle
    - \frac{\kappa}{2}\| \mathbf{v} - \mathbf{v}_B \|^2.
\]

Where:
- first term: beliefs $B$ amplify policy-gradient ascent on $\Phi$
- second term: $\mathbf{v}_B$ is the belief-implied optimal policy

This is the field-theoretic embodiment of:
\[
\text{Best policy} = \text{argmax predicted welfare}.
\]

% -------------------
\subsubsection*{5. Entropy-Field Term}

Entropy penalizes sharp, unstable, or overly complex structures:
\[
\mathcal{L}_{\text{entropy}}
=
- \eta\, |\nabla S|^2 
- \xi\, S \log S.
\]

The second term is Shannon entropy.

% ------------------------------------------------------------
\subsection*{E. Euler-Lagrange Field Equations}

We now vary the action with respect to each field.

\begin{lemma}[Euler-Lagrange Equations for RSVP-Futarchy Fields]
The fields satisfy:
\[
\frac{\partial \mathcal{L}}{\partial X}
-
\partial_\mu
\left(
\frac{\partial \mathcal{L}}{\partial(\partial_\mu X)}
\right)
= 0,
\quad
X \in \{\Phi, \mathbf{v}, S, B\}.
\]
\end{lemma}

Solving yields:

\subsubsection*{Value Field Equation}
\[
\partial_t^2 \Phi
= 
\lambda\, \frac{\delta}{\delta \Phi}\|T\|^2
+ \gamma\, \nabla \cdot (B \mathbf{v}).
\]

\subsubsection*{Policy Flow Equation}
\[
\partial_t^2 \mathbf{v}
=
\lambda\, \frac{\delta}{\delta \mathbf{v}}\|T\|^2
+
\gamma\, B\, \nabla \Phi
-
\kappa( \mathbf{v} - \mathbf{v}_B ).
\]

\subsubsection*{Belief Field Equation}
\[
\partial_t^2 B
=
\alpha\, \Delta B
+ \beta\, S
+ \gamma\, \langle \nabla \Phi, \mathbf{v} \rangle.
\]

\subsubsection*{Entropy Field Equation}
\[
\partial_t^2 S
=
\lambda \frac{\delta}{\delta S}\|T\|^2
+ \beta B
+ 2\eta \Delta S
+ \xi(1+\log S).
\]

% ------------------------------------------------------------
\subsection*{F. Main Result}

\begin{theorem}[RSVP-Futarchy Equivalence Principle]
A field configuration
\[
(\Phi^\ast, \mathbf{v}^\ast, S^\ast, B^\ast)
\]
is a stationary solution of the futarchy-Lagrangian if and only if:
\begin{enumerate}
\item it satisfies the economic equilibrium condition
\[
T(\Phi^\ast,\mathbf{v}^\ast,S^\ast)=0,
\]
\item the policy maximizes predicted welfare
\[
\mathbf{v}^\ast = \arg\max_{\mathbf{v}}
\ \mathbb{E}_{B^\ast}\!\left[\Phi^\ast\right],
\]
\item beliefs are predictive-accuracy stationary points,
\item and entropy is minimized subject to the above.
\end{enumerate}
\end{theorem}

\begin{proof}
Stationarity of the action produces the Euler-Lagrange equations above.
The $\lambda$-term enforces equilibrium. The $\kappa$- and $\gamma$-terms
impose the futarchy decision rule. The remaining terms enforce belief
accuracy and entropy minimization. Thus all conditions follow.
\end{proof}

% ============================================================
% Appendix: Hamiltonian Formulation of RSVP-Futarchy Fields
% ============================================================

\section*{Appendix: Hamiltonian Formulation of the RSVP-Futarchy System}

In this appendix we derive the canonical Hamiltonian formulation of
the RSVP-Futarchy field theory introduced in the previous appendix.
We work with the same field content:
\[
\Phi(\mathbf{x},t),\quad
\mathbf{v}(\mathbf{x},t),\quad
S(\mathbf{x},t),\quad
B(\mathbf{x},t),
\]
and the Lagrangian density
\[
\mathcal{L}
=
\mathcal{L}_{\text{dyn}}
+
\mathcal{L}_{\text{equil}}
+
\mathcal{L}_{\text{belief}}
+
\mathcal{L}_{\text{futarchy}}
+
\mathcal{L}_{\text{entropy}},
\]
where
\begin{align}
\mathcal{L}_{\text{dyn}}
&=
\frac{1}{2}\Big(
    |\partial_t \Phi|^2
  + \|\partial_t \mathbf{v}\|^2
  + |\partial_t S|^2
  + |\partial_t B|^2
\Big), \\[4pt]
\mathcal{L}_{\text{equil}}
&=
- \frac{\lambda}{2}\,
    \|T(\Phi, \mathbf{v}, S)\|^2, \\[4pt]
\mathcal{L}_{\text{belief}}
&=
- \frac{\alpha}{2}
    \|\nabla B\|^2
    - \beta\, B\, S, \\[4pt]
\mathcal{L}_{\text{futarchy}}
&=
\gamma\, B\,\langle \nabla \Phi,\ \mathbf{v} \rangle
    - \frac{\kappa}{2}\| \mathbf{v} - \mathbf{v}_B \|^2, \\[4pt]
\mathcal{L}_{\text{entropy}}
&=
- \eta\, |\nabla S|^2 
- \xi\, S \log S.
\end{align}

% ------------------------------------------------------------
\subsection*{A. Canonical Momenta and Phase Space}

We define the conjugate momenta in the standard way:
\[
\Pi_X(\mathbf{x},t)
=
\frac{\partial \mathcal{L}}{\partial (\partial_t X(\mathbf{x},t))},
\qquad
X \in \{\Phi,\mathbf{v},S,B\}.
\]

Since only $\mathcal{L}_{\text{dyn}}$ depends on time derivatives, the
canonical momenta are:
\begin{align}
\Pi_\Phi &= \frac{\partial \mathcal{L}}{\partial(\partial_t \Phi)}
        = \partial_t \Phi, \\[4pt]
\bm{\Pi}_{\mathbf{v}} &= 
\frac{\partial \mathcal{L}}{\partial(\partial_t \mathbf{v})}
        = \partial_t \mathbf{v}, \\[4pt]
\Pi_S &= \frac{\partial \mathcal{L}}{\partial(\partial_t S)}
        = \partial_t S, \\[4pt]
\Pi_B &= \frac{\partial \mathcal{L}}{\partial(\partial_t B)}
        = \partial_t B.
\end{align}

\begin{definition}[RSVP--Futarchy Phase Space]
The canonical phase space consists of the fields and their conjugate
momenta:
\[
\big( \Phi, \Pi_\Phi;\ 
      \mathbf{v}, \bm{\Pi}_{\mathbf{v}};\ 
      S, \Pi_S;\ 
      B, \Pi_B
\big).
\]
\end{definition}

In this formulation, the lamphrodynamic evolution is a Hamiltonian flow
on this infinite-dimensional phase space.

% ------------------------------------------------------------
\subsection*{B. Hamiltonian Density}

The Hamiltonian density is defined as
\[
\mathcal{H}
= 
\Pi_\Phi\, \partial_t \Phi
+
\bm{\Pi}_{\mathbf{v}}\!\cdot \partial_t \mathbf{v}
+
\Pi_S\, \partial_t S
+
\Pi_B\, \partial_t B
-
\mathcal{L}.
\]

Substituting $\partial_t \Phi = \Pi_\Phi$, etc., we obtain
\begin{equation}
\Pi_\Phi\, \partial_t \Phi
+
\bm{\Pi}_{\mathbf{v}}\!\cdot \partial_t \mathbf{v}
+
\Pi_S\, \partial_t S
+
\Pi_B\, \partial_t B
=
|\Pi_\Phi|^2
+
\|\bm{\Pi}_{\mathbf{v}}\|^2
+
|\Pi_S|^2
+
|\Pi_B|^2
= 2\, \mathcal{L}_{\text{dyn}}.
\end{equation}

Therefore,
\begin{align}
\mathcal{H}
&=
2\mathcal{L}_{\text{dyn}} 
-
\big(
\mathcal{L}_{\text{dyn}}
+
\mathcal{L}_{\text{equil}}
+
\mathcal{L}_{\text{belief}}
+
\mathcal{L}_{\text{futarchy}}
+
\mathcal{L}_{\text{entropy}}
\big) \\[4pt]
&=
\mathcal{L}_{\text{dyn}}
-
\mathcal{L}_{\text{equil}}
-
\mathcal{L}_{\text{belief}}
-
\mathcal{L}_{\text{futarchy}}
-
\mathcal{L}_{\text{entropy}}.
\end{align}

Using the explicit forms of each term:
\begin{align}
\mathcal{H}
&=
\frac{1}{2}\Big(
    |\Pi_\Phi|^2
  + \|\bm{\Pi}_{\mathbf{v}}\|^2
  + |\Pi_S|^2
  + |\Pi_B|^2
\Big)
+ \frac{\lambda}{2}\,\|T(\Phi,\mathbf{v},S)\|^2 \\[4pt]
&\quad
+ \frac{\alpha}{2}\|\nabla B\|^2
+ \beta\, B S \\[4pt]
&\quad
- \gamma\, B\,\langle \nabla \Phi,\ \mathbf{v} \rangle
+ \frac{\kappa}{2}\|\mathbf{v} - \mathbf{v}_B\|^2 \\[4pt]
&\quad
+ \eta\,|\nabla S|^2
+ \xi\, S \log S.
\end{align}

\begin{definition}[RSVP--Futarchy Hamiltonian]
The Hamiltonian functional is
\[
H[\Phi,\mathbf{v},S,B;\Pi_\Phi,\bm{\Pi}_{\mathbf{v}},\Pi_S,\Pi_B]
=
\int_{\mathcal{X}}
\mathcal{H}\, d\mu.
\]
\end{definition}

This $H$ plays the role of a generalized free energy, combining:
kinetic contributions from the momenta, equilibrium violations, belief
curvature, futarchy couplings, and entropy penalties.

% ------------------------------------------------------------
\subsection*{C. Hamilton’s Equations}

The canonical Hamilton equations take the form
\begin{align}
\partial_t X
&= \frac{\delta H}{\delta \Pi_X}, \\[4pt]
\partial_t \Pi_X
&= - \frac{\delta H}{\delta X},
\end{align}
for each field $X \in \{\Phi,\mathbf{v},S,B\}$.

\subsubsection*{1. Value Field}
\begin{align}
\partial_t \Phi
&= \frac{\delta H}{\delta \Pi_\Phi}
= \Pi_\Phi, \\[4pt]
\partial_t \Pi_\Phi
&= -\frac{\delta H}{\delta \Phi}.
\end{align}
The second equation encodes the effects of:
\begin{itemize}
    \item equilibrium curvature via $\lambda\,\|T\|^2$,
    \item futarchy coupling via $-\gamma\, B\,\langle \nabla (\cdot),\mathbf{v}\rangle$.
\end{itemize}

\subsubsection*{2. Policy Field}
\begin{align}
\partial_t \mathbf{v}
&= \frac{\delta H}{\delta \bm{\Pi}_{\mathbf{v}}}
= \bm{\Pi}_{\mathbf{v}}, \\[4pt]
\partial_t \bm{\Pi}_{\mathbf{v}}
&= -\frac{\delta H}{\delta \mathbf{v}}.
\end{align}
The functional derivative includes:
\begin{itemize}
    \item equilibrium penalties,
    \item a futarchy term $+\gamma B \nabla \Phi$,
    \item a restoring term $-\kappa (\mathbf{v}-\mathbf{v}_B)$.
\end{itemize}

\subsubsection*{3. Belief Field}
\begin{align}
\partial_t B
&= \frac{\delta H}{\delta \Pi_B}
= \Pi_B, \\[4pt]
\partial_t \Pi_B
&= -\frac{\delta H}{\delta B}.
\end{align}
The potential term in $H$ gives:
\[
\frac{\delta H}{\delta B}
=
\beta S
- \gamma\,\langle \nabla \Phi,\mathbf{v} \rangle
- \alpha \Delta B,
\]
so
\[
\partial_t^2 B
=
\alpha\,\Delta B
+ \beta S
+ \gamma\,\langle \nabla \Phi,\mathbf{v} \rangle,
\]
matching the Euler--Lagrange equation.

\subsubsection*{4. Entropy Field}
\begin{align}
\partial_t S
&= \frac{\delta H}{\delta \Pi_S}
= \Pi_S, \\[4pt]
\partial_t \Pi_S
&= -\frac{\delta H}{\delta S}.
\end{align}
The potential includes:
\[
\frac{\delta H}{\delta S}
=
\lambda\,\frac{\delta}{\delta S}\|T\|^2
+ \beta B
+ 2\eta \Delta S
+ \xi(1+\log S),
\]
giving the same entropy evolution as in the Lagrangian formulation.

\begin{proposition}[Equivalence with Euler--Lagrange Dynamics]
The Hamiltonian equations derived above are equivalent to the
Euler--Lagrange field equations obtained from the RSVP--Futarchy
Lagrangian.
\end{proposition}

\begin{proof}
Direct verification: for each field, the Hamiltonian evolution
reproduces the second-order time equations obtained from the
Euler--Lagrange formalism, given the identification of momenta
$\Pi_X = \partial_t X$.
\end{proof}

% ------------------------------------------------------------
\subsection*{D. Conceptual Interpretation}

The Hamiltonian formulation emphasizes that the RSVP--Futarchy system
is a \emph{phase-space dynamical system} in which:
\begin{itemize}
    \item the fields $(\Phi,\mathbf{v},S,B)$ encode welfare, policy,
    entropy, and collective beliefs;
    \item the conjugate momenta represent the ``lamphrodynamic
    velocities'' of these structures;
    \item the Hamiltonian $H$ acts as a generalized free energy,
    combining equilibrium constraints, prediction markets, and
    entropy penalties;
    \item futarchy couplings appear as cross-terms in the potential
    energy that direct policy flows along those trajectories that
    markets predict will maximize welfare.
\end{itemize}

Stationary points of $H$ correspond to fully self-consistent economic
equilibria in which values, beliefs, policies, and entropy are jointly
aligned. The Hamiltonian view thus complements the Lagrangian
field-theoretic picture and provides a natural foundation for canonical
quantization or stochastic extensions of predictive market dynamics.

% ============================================================
% Appendix: Poisson Brackets and Symplectic Structure
%          for the RSVP--Futarchy Hamiltonian System
% ============================================================

\section*{Appendix: Poisson Brackets and Symplectic Structure\\
for RSVP--Futarchy Fields}

In this appendix we endow the RSVP--Futarchy Hamiltonian system with a
canonical symplectic structure and the associated Poisson bracket
formalism. This makes explicit that the lamphrodynamic evolution of
welfare, policy, entropy, and belief fields is a Hamiltonian flow on an
infinite-dimensional phase space.

We continue with the field content and canonical momenta introduced in
the Hamiltonian appendix:
\[
\Phi(\mathbf{x},t),\quad
\mathbf{v}(\mathbf{x},t),\quad
S(\mathbf{x},t),\quad
B(\mathbf{x},t),
\]
with conjugate momenta
\[
\Pi_\Phi(\mathbf{x},t),\quad
\bm{\Pi}_{\mathbf{v}}(\mathbf{x},t),\quad
\Pi_S(\mathbf{x},t),\quad
\Pi_B(\mathbf{x},t),
\]
and Hamiltonian functional
\[
H[\Phi,\mathbf{v},S,B;\Pi_\Phi,\bm{\Pi}_{\mathbf{v}},\Pi_S,\Pi_B]
=
\int_{\mathcal{X}} \mathcal{H} \, d\mu.
\]

% ------------------------------------------------------------
\subsection*{A. Canonical Poisson Brackets}

We treat the fields and their conjugate momenta as coordinates on an
infinite-dimensional phase space and define the canonical Poisson
brackets in the standard functional sense.

Let $F$ and $G$ be functionals of the fields and momenta. The
\emph{canonical Poisson bracket} is defined as
\begin{equation}
\{F,G\}
=
\int_{\mathcal{X}}
\left[
    \frac{\delta F}{\delta \Phi(\mathbf{x})}
    \,\frac{\delta G}{\delta \Pi_\Phi(\mathbf{x})}
  - \frac{\delta F}{\delta \Pi_\Phi(\mathbf{x})}
    \,\frac{\delta G}{\delta \Phi(\mathbf{x})}
\right] d\mu(\mathbf{x})
+ \cdots,
\end{equation}
with analogous contributions from the other field--momentum pairs:
\[
(\mathbf{v},\bm{\Pi}_{\mathbf{v}}),\quad (S,\Pi_S),\quad (B,\Pi_B).
\]

More explicitly,
\begin{align}
\{F,G\}
=&
\int_{\mathcal{X}}
\left[
    \frac{\delta F}{\delta \Phi}
    \,\frac{\delta G}{\delta \Pi_\Phi}
  - \frac{\delta F}{\delta \Pi_\Phi}
    \,\frac{\delta G}{\delta \Phi}
\right] d\mu
\\[4pt]
&+
\int_{\mathcal{X}}
\left[
    \frac{\delta F}{\delta v^i}
    \,\frac{\delta G}{\delta \Pi_{v^i}}
  - \frac{\delta F}{\delta \Pi_{v^i}}
    \,\frac{\delta G}{\delta v^i}
\right] d\mu
\\[4pt]
&+
\int_{\mathcal{X}}
\left[
    \frac{\delta F}{\delta S}
    \,\frac{\delta G}{\delta \Pi_S}
  - \frac{\delta F}{\delta \Pi_S}
    \,\frac{\delta G}{\delta S}
\right] d\mu
\\[4pt]
&+
\int_{\mathcal{X}}
\left[
    \frac{\delta F}{\delta B}
    \,\frac{\delta G}{\delta \Pi_B}
  - \frac{\delta F}{\delta \Pi_B}
    \,\frac{\delta G}{\delta B}
\right] d\mu,
\end{align}
where repeated spatial indices $i$ are summed and we suppress explicit
$\mathbf{x}$-dependence for readability.

\begin{definition}[Canonical Equal-Time Poisson Brackets]
The fundamental brackets between the fields and their conjugate
momenta at equal time are:
\begin{align}
\{\Phi(\mathbf{x}), \Pi_\Phi(\mathbf{y})\}
    &= \delta(\mathbf{x}-\mathbf{y}), &
\{\Phi(\mathbf{x}), \Phi(\mathbf{y})\} &= 0, &
\{\Pi_\Phi(\mathbf{x}), \Pi_\Phi(\mathbf{y})\} &= 0,
\\[4pt]
\{v^i(\mathbf{x}), \Pi_{v^j}(\mathbf{y})\}
    &= \delta^i{}_j\,\delta(\mathbf{x}-\mathbf{y}), &
\{v^i(\mathbf{x}), v^j(\mathbf{y})\} &= 0, &
\{\Pi_{v^i}(\mathbf{x}), \Pi_{v^j}(\mathbf{y})\} &= 0,
\\[4pt]
\{S(\mathbf{x}), \Pi_S(\mathbf{y})\}
    &= \delta(\mathbf{x}-\mathbf{y}), &
\{S(\mathbf{x}), S(\mathbf{y})\} &= 0, &
\{\Pi_S(\mathbf{x}), \Pi_S(\mathbf{y})\} &= 0,
\\[4pt]
\{B(\mathbf{x}), \Pi_B(\mathbf{y})\}
    &= \delta(\mathbf{x}-\mathbf{y}), &
\{B(\mathbf{x}), B(\mathbf{y})\} &= 0, &
\{\Pi_B(\mathbf{x}), \Pi_B(\mathbf{y})\} &= 0,
\end{align}
with all cross-brackets between distinct fields vanishing.
\end{definition}

These Poisson brackets satisfy the usual antisymmetry, bilinearity, 
Leibniz rule, and Jacobi identity, and thus equip the function space
with a Lie-algebra structure.

% ------------------------------------------------------------
\subsection*{B. Symplectic 2-Form}

The canonical brackets arise from a symplectic 2-form on phase space.

\begin{definition}[RSVP--Futarchy Symplectic Form]
The symplectic form on the phase space of fields is
\begin{equation}
\Omega
=
\int_{\mathcal{X}}
\Big(
    \delta \Pi_\Phi \wedge \delta \Phi
  + \delta \bm{\Pi}_{\mathbf{v}} \wedge \delta \mathbf{v}
  + \delta \Pi_S \wedge \delta S
  + \delta \Pi_B \wedge \delta B
\Big)\, d\mu.
\end{equation}
\end{definition}

Here $\delta$ denotes the exterior derivative on the infinite-dimensional 
field space, and $\wedge$ is the wedge product. Non-degeneracy and
closedness ($d\Omega=0$) ensure that $(\text{phase space},\Omega)$ is a
symplectic manifold in the functional sense.

\begin{lemma}[Poisson Brackets from $\Omega$]
The Poisson bracket of functionals $F$ and $G$ can be written as
\[
\{F,G\} = \Omega(X_F,X_G),
\]
where $X_F$ is the Hamiltonian vector field associated with $F$, defined
implicitly by
\[
\iota_{X_F} \Omega = \delta F.
\]
\end{lemma}

\begin{proof}
Standard symplectic geometry argument: the definition of $X_F$ via
$\iota_{X_F}\Omega=\delta F$ implies that for any $G$,
\[
\{F,G\} := \Omega(X_F,X_G)
= \delta F(X_G),
\]
which reproduces the functional expression given earlier.
\end{proof}

Thus the canonical brackets and symplectic form are equivalent ways of
encoding the same geometric structure.

% ------------------------------------------------------------
\subsection*{C. Hamiltonian Flow as Lamphrodynamic Evolution}

With the symplectic form and Poisson bracket defined, the evolution of
any functional $F$ of the fields and momenta under the Hamiltonian
\[
H = \int_{\mathcal{X}}\mathcal{H}\,d\mu
\]
is given by
\[
\frac{dF}{dt} = \{F,H\}.
\]

In particular, for the fields themselves:
\begin{align}
\partial_t \Phi &= \{\Phi,H\} = \frac{\delta H}{\delta \Pi_\Phi}, &
\partial_t \Pi_\Phi &= \{\Pi_\Phi,H\} = -\frac{\delta H}{\delta \Phi},\\[4pt]
\partial_t \mathbf{v} &= \{\mathbf{v},H\} = \frac{\delta H}{\delta \bm{\Pi}_{\mathbf{v}}}, &
\partial_t \bm{\Pi}_{\mathbf{v}} &= \{\bm{\Pi}_{\mathbf{v}},H\} = -\frac{\delta H}{\delta \mathbf{v}},\\[4pt]
\partial_t S &= \{S,H\} = \frac{\delta H}{\delta \Pi_S}, &
\partial_t \Pi_S &= \{\Pi_S,H\} = -\frac{\delta H}{\delta S},\\[4pt]
\partial_t B &= \{B,H\} = \frac{\delta H}{\delta \Pi_B}, &
\partial_t \Pi_B &= \{\Pi_B,H\} = -\frac{\delta H}{\delta B}.
\end{align}

\begin{proposition}[Lamphrodynamic Flow as Hamiltonian Flow]
The lamphrodynamic evolution derived from the Lagrangian formulation is 
identical to the Hamiltonian flow generated by $H$ with respect to the
canonical Poisson brackets.
\end{proposition}

\begin{proof}
Substituting the Hamiltonian density $\mathcal{H}$ into the functional
derivatives above reproduces the Euler--Lagrange equations for each
field, given the identification $\Pi_X = \partial_t X$. This is the
standard equivalence between the Lagrangian and Hamiltonian pictures.
\end{proof}

Thus the RSVP--Futarchy system is a fully symplectic dynamical system,
with lamphrodynamic evolution interpreted as a Hamiltonian flow on the
space of welfare, policy, entropy, and belief fields.

% ------------------------------------------------------------
\subsection*{D. Remarks on Constraints and Futarchy Couplings}

In the present construction we have taken the field variables and
momenta as unconstrained canonical coordinates. In more refined models:
\begin{itemize}
    \item the equilibrium condition $T(\Phi,\mathbf{v},S)=0$ can be
    implemented as a \emph{constraint surface} in phase space;
    \item the futarchy relationship 
    $\mathbf{v} \approx \mathbf{v}_B(\Phi,B)$ may introduce additional
    (possibly second-class) constraints;
    \item value-gauge symmetries (redefining $\Phi$ by additive
    constants or linear transformations) may appear as first-class
    constraints generating gauge orbits.
\end{itemize}

A full Dirac constraint analysis would then identify reduced phase
spaces and modified (Dirac) brackets. For the purposes of this essay,
however, the canonical Poisson structure suffices to demonstrate that
the RSVP--Futarchy fields form a legitimate Hamiltonian field theory,
ready for possible quantization or stochastic deformation.

% ============================================================
% Appendix: Path Integrals and Functional Integrals
%           for RSVP–Futarchy Field Theory
% ============================================================

\section*{Appendix: Path and Functional Integrals\\
for the RSVP--Futarchy Field Theory}

This appendix provides a functional-integral formulation of the
RSVP–Futarchy system. In this picture, welfare fields, policy fields,
entropy fields, and belief fields are integrated over all possible
configurations weighted by the futarchy Lagrangian. The result is a
field-theoretic interpretation of prediction markets as information-
aggregating functional integrators.

% ------------------------------------------------------------
\subsection*{A. The Action and Functional Measure}

We begin with the Lagrangian density
\[
\mathcal{L}
=
\mathcal{L}_{\text{dyn}}
+
\mathcal{L}_{\text{equil}}
+
\mathcal{L}_{\text{belief}}
+
\mathcal{L}_{\text{futarchy}}
+
\mathcal{L}_{\text{entropy}},
\]
with explicit forms given in the earlier appendix.

The action is
\[
S[\Phi,\mathbf{v},S,B]
=
\int dt \int_{\mathcal{X}}
\mathcal{L}(\Phi,\mathbf{v},S,B)\, d\mu.
\]

The formal functional measure is
\[
\mathcal{D}\Phi
\,\mathcal{D}\mathbf{v}
\,\mathcal{D}S
\,\mathcal{D}B,
\]
which integrates over all histories of the fields.

% ------------------------------------------------------------
\subsection*{B. The RSVP--Futarchy Generating Functional}

The central object is the generating functional:
\begin{equation}
\boxed{
Z[J_\Phi, J_{\mathbf{v}}, J_S, J_B]
=
\int
\mathcal{D}\Phi\,
\mathcal{D}\mathbf{v}\,
\mathcal{D}S\,
\mathcal{D}B\;
\exp\!\left[
    - S[\Phi,\mathbf{v},S,B]
    + \int J \cdot X
\right].
}
\end{equation}

Here:
\[
J\cdot X
=
\int dt \int_{\mathcal{X}}
\big( 
J_\Phi \Phi
+ J_{\mathbf{v}}\cdot\mathbf{v}
+ J_S S
+ J_B B
\big)\, d\mu.
\]

We adopt the Euclidean convention (\(\exp[-S]\)) for stability and
interpretability.

Correlation functions follow by functional differentiation:
\[
\left\langle X_1(\mathbf{x}_1,t_1)\cdots X_n(\mathbf{x}_n,t_n)\right\rangle
=
\frac{\delta^n Z}{\delta J_1 \cdots \delta J_n}
\Bigg|_{J=0}.
\]

% ------------------------------------------------------------
\subsection*{C. Integrating Out the Belief Field $B$}

A central feature of futarchy is that predictive markets aggregate
information into a single, coherent belief distribution. This property
emerges naturally from the functional integral.

We isolate all terms involving $B$ in the Euclidean action:
\begin{align}
S_B
&=
\int dt \int_{\mathcal{X}}
\left[
\frac{1}{2}|\partial_t B|^2
+ \frac{\alpha}{2}\|\nabla B\|^2
+ \beta B S
- \gamma B \langle \nabla \Phi,\mathbf{v}\rangle
\right] d\mu.
\end{align}

This is a quadratic functional of $B$ with linear source:
\[
J_B^{\text{eff}}
=
\beta S
- \gamma \langle \nabla \Phi,\mathbf{v} \rangle.
\]

Thus the belief integral is Gaussian:
\begin{equation}
\int \mathcal{D}B\;
\exp\!\left[ -S_B\right]
=
\big(\det \mathcal{O}_B\big)^{-1/2}
\exp\!\left[
    \frac{1}{2}
    \int J_B^{\text{eff}}
    \mathcal{O}_B^{-1}
    J_B^{\text{eff}}
\right],
\end{equation}
where
\[
\mathcal{O}_B
=
- \partial_t^2
+ \alpha \Delta
\]
is the belief-field kinetic operator.

\begin{theorem}[Futarchy as Optimal Information Aggregation]
Integrating out the belief field produces a nonlocal effective action
that encodes optimal prediction aggregation:
\[
S_{\text{eff}}
= S[\Phi,\mathbf{v},S]
- \frac{1}{2}
    \int
    J_B^{\text{eff}}
    \mathcal{O}_B^{-1}
    J_B^{\text{eff}}.
\]
\end{theorem}

\begin{proof}
Gaussian functional integration. The inverse operator $\mathcal{O}_B^{-1}$
is the Green’s function that aggregates information from all points in
spacetime.
\end{proof}

This term is explicitly nonlocal:
\[
\mathcal{O}_B^{-1}(\mathbf{x},t;\mathbf{y},t')
\sim
\text{propagator connecting all prediction sources}.
\]

Economically, it captures how speculative markets pool information,
compare predictions globally, and feed back into policy flows.

% ------------------------------------------------------------
\subsection*{D. Effective Action After Integrating Out $B$}

The effective action is:
\[
S_{\text{eff}}
=
\int dt \int_{\mathcal{X}}
\left[
\mathcal{L}_{\text{dyn}}
+
\mathcal{L}_{\text{equil}}
+
\mathcal{L}_{\text{entropy}}
\right] d\mu
+
S_{\text{belief-induced}}[\Phi,\mathbf{v},S],
\]
with
\[
S_{\text{belief-induced}}
=
- \frac{1}{2}
\int dt\,dt'\,\int_{\mathcal{X}}d\mu\int_{\mathcal{X}}d\mu'
\,
J_B^{\text{eff}}(x)
\,
\mathcal{O}_B^{-1}(x,x')
\,
J_B^{\text{eff}}(x').
\]

This contains cross-terms of the form:
\begin{align}
&\phantom{{}={}}
\frac{1}{2}\gamma^2
\langle \nabla\Phi,\mathbf{v}\rangle(x)
\,\mathcal{O}_B^{-1}(x,x')\,
\langle \nabla\Phi,\mathbf{v}\rangle(x'),\\
&\phantom{{}={}}
- \gamma \beta\,
\langle \nabla\Phi,\mathbf{v}\rangle(x)
\,\mathcal{O}_B^{-1}(x,x')\,
S(x'),\\
&\phantom{{}={}}
\frac{1}{2}\beta^2\,
S(x)\,\mathcal{O}_B^{-1}(x,x')\,S(x').
\end{align}

\begin{definition}[Futarchy Kernel]
Define the kernel
\[
K(x,x') := \mathcal{O}_B^{-1}(x,x')
\]
as the belief-propagation operator.
\end{definition}

Then
\[
S_{\text{belief-induced}}
=
- \frac{1}{2}
\iint
J_B^{\text{eff}}(x)
\,K(x,x')\,
J_B^{\text{eff}}(x')
\,dx\,dx'.
\]

This confirms that the market-belief field creates long-range
interactions across the economic manifold, exactly as futarchy
predicts.

% ------------------------------------------------------------
\subsection*{E. Probability Amplitude for Policy Paths}

The probability amplitude for a policy trajectory $\mathbf{v}(t)$ is:
\[
\mathcal{P}[\mathbf{v}]
\propto
\int
\mathcal{D}\Phi\,\mathcal{D}S
\;
\exp[-S_{\text{eff}}(\Phi,\mathbf{v},S)].
\]

The most probable trajectories minimize the effective action:
\[
\frac{\delta S_{\text{eff}}}{\delta \mathbf{v}} = 0,
\]
which yields the futarchy-adjusted Euler–Lagrange equations.

This is the field-theoretic proof that:

> *Prediction markets select the policies that maximize expected value.*

% ------------------------------------------------------------
\subsection*{F. Functional Integral Interpretation of Futarchy}

This framework yields the following interpretation:

\begin{itemize}
    \item Welfare potential \(\Phi\) is what society votes on.  
    \item Policy field \(\mathbf{v}\) is what society chooses (or what the system evolves toward).  
    \item The belief field \(B\) is integrated out, leaving behind a nonlocal kernel \(K = \mathcal{O}^{-1}\) that aggregates all predictive information.  
    \item Policy is chosen to minimize the effective action, i.e., to maximize expected welfare under market-generated probabilistic beliefs.
\end{itemize}

Thus futarchy naturally emerges as:

\begin{quote}
\textbf{Value selection via variational optimization and belief aggregation via functional integration.}
\end{quote}

This gives a clean mathematical foundation for the slogan:
\[
\text{Vote on values; bet on beliefs.}
\]

% ============================================================
% Appendix: Renormalization of the Belief Field
%           and RG Flow for Futarchy Couplings
% ============================================================

\section*{Appendix: Renormalization of the Belief Field\\
and Scale-Dependent Futarchy Couplings}

In this appendix we analyze the RSVP--Futarchy field theory from a
Wilsonian renormalization-group (RG) viewpoint. The central idea is that
the belief field $B$ acts as a scale-dependent information aggregator:
short-scale fluctuations represent noisy, local bets, while long-scale
modes encode robust, macro-level predictions. By integrating out
high-frequency components of $B$, we obtain an effective theory with
scale-dependent couplings, providing a renormalization interpretation of
futarchy.

% ------------------------------------------------------------
\subsection*{A. Momentum-Shell Decomposition of the Belief Field}

For simplicity, we work in a flat spatial manifold with coordinates
$\mathbf{x}\in\mathbb{R}^d$ and use a Fourier representation. The
extension to curved $\mathcal{X}$ uses eigenmodes of the Laplacian.

Let
\[
B(\mathbf{x},t)
=
\int \frac{d^d k}{(2\pi)^d}\;
e^{i\mathbf{k}\cdot\mathbf{x}}\,
\tilde{B}(\mathbf{k},t).
\]

We introduce a UV cutoff $\Lambda$ and decompose modes into:
\begin{align}
\tilde{B}_{<}(\mathbf{k},t) &: \|\mathbf{k}\| \le \Lambda/b,\\
\tilde{B}_{>}(\mathbf{k},t) &: \Lambda/b < \|\mathbf{k}\| \le \Lambda,
\end{align}
for some rescaling factor $b>1$. Correspondingly,
\[
B = B_{<} + B_{>}.
\]

\begin{definition}[Belief Field Mode Split]
The low-frequency component $B_{<}$ encodes large-scale, aggregated 
predictive structure, while $B_{>}$ encodes short-scale, local
fluctuations in beliefs.
\end{definition}

The Wilsonian RG step integrates out $B_{>}$ to obtain an effective
action $S_{\text{eff}}[\Phi,\mathbf{v},S,B_{<};\ell]$ at scale
$\ell = \log b$.

% ------------------------------------------------------------
\subsection*{B. Belief Sector of the Action}

Recall the $B$-dependent part of the Euclidean action:
\[
S_B[B;\Phi,\mathbf{v},S]
=
\int dt \int_{\mathcal{X}}
\left[
\frac{1}{2}|\partial_t B|^2
+ \frac{\alpha}{2}\|\nabla B\|^2
+ \beta B S
- \gamma B \langle \nabla \Phi,\mathbf{v}\rangle
\right] d\mu.
\]

In Fourier space (suppressing time and indices),
\[
S_B
=
\frac{1}{2}
\int_{\|\mathbf{k}\|\le\Lambda}
\tilde{B}(-\mathbf{k})\,
\mathcal{O}_B(\mathbf{k})\,
\tilde{B}(\mathbf{k})\, d^d k
+
\int_{\|\mathbf{k}\|\le\Lambda}
\tilde{B}(-\mathbf{k})\,\tilde{J}_B(\mathbf{k})\, d^d k,
\]
with
\[
\mathcal{O}_B(\mathbf{k})
=
\omega^2 + \alpha \|\mathbf{k}\|^2, 
\qquad
\tilde{J}_B(\mathbf{k})
=
\widetilde{\beta S}
    - \widetilde{\gamma \langle\nabla\Phi,\mathbf{v}\rangle}.
\]

We now perform a Wilsonian integration of $B_{>}$.

% ------------------------------------------------------------
\subsection*{C. Wilsonian Integration of High-Frequency Belief Modes}

Decompose
\[
\tilde{B}(\mathbf{k})
=
\tilde{B}_{<}(\mathbf{k})\,\theta(\Lambda/b - \|\mathbf{k}\|)
+
\tilde{B}_{>}(\mathbf{k})\,\theta(\|\mathbf{k}\|-\Lambda/b)\theta(\Lambda-\|\mathbf{k}\|).
\]

The partition function becomes
\[
Z
=
\int \mathcal{D}\Phi\,\mathcal{D}\mathbf{v}\,\mathcal{D}S
\int \mathcal{D}B_{<}\,\mathcal{D}B_{>}
\;
e^{-S[\Phi,\mathbf{v},S,B_{<}+B_{>}]}.
\]

The Wilsonian step defines an effective action at scale $\Lambda/b$ via
\[
e^{-S_{\text{eff}}[\Phi,\mathbf{v},S,B_{<};\Lambda/b]}
=
\int \mathcal{D}B_{>}\;
e^{-S_B[B_{<}+B_{>};\Phi,\mathbf{v},S]}
\; e^{-S_{\text{rest}}[\Phi,\mathbf{v},S]},
\]
where $S_{\text{rest}}$ contains all other fields’ contributions.

Because $S_B$ is quadratic in $B$, the integral over $B_{>}$ is again
Gaussian, yielding a correction term to the effective action for $B_{<}$:
\begin{equation}
S_{\text{eff}}[\Phi,\mathbf{v},S,B_{<}]
=
S_B[B_{<};\Phi,\mathbf{v},S]
+ \Delta S_{\Lambda\to\Lambda/b}[\Phi,\mathbf{v},S]
+ S_{\text{rest}}[\Phi,\mathbf{v},S].
\end{equation}

The correction term $\Delta S$ arises from the contraction of the source
$\tilde{J}_B$ with the propagator over the high-momentum shell:
\begin{align}
\Delta S_{\Lambda\to\Lambda/b}
&=
- \frac{1}{2}
\int_{\Lambda/b < \|\mathbf{k}\| \le \Lambda}
\frac{|\tilde{J}_B(\mathbf{k})|^2}{\mathcal{O}_B(\mathbf{k})}
\, d^d k.
\end{align}

This correction renormalizes the couplings that multiply $S$ and
$\langle\nabla\Phi,\mathbf{v}\rangle$ in the low-energy theory.

% ------------------------------------------------------------
\subsection*{D. Running Couplings and Beta Functions}

At leading order, we can treat $\tilde{J}_B$ as small and expand
$\Delta S$ to extract the induced running of parameters. The belief
sector couplings $(\alpha,\beta,\gamma)$ now become scale-dependent:
\[
\alpha \to \alpha(\ell),\quad
\beta \to \beta(\ell),\quad
\gamma \to \gamma(\ell),
\]
where $\ell = \log(\Lambda/\Lambda_{\text{eff}})$.

\begin{definition}[RG Flow of Futarchy Couplings]
We define beta functions by
\[
\frac{d\alpha}{d\ell} = \beta_\alpha(\alpha,\beta,\gamma,\dots),
\quad
\frac{d\beta}{d\ell} = \beta_\beta(\alpha,\beta,\gamma,\dots),
\quad
\frac{d\gamma}{d\ell} = \beta_\gamma(\alpha,\beta,\gamma,\dots).
\]
\end{definition}

Qualitatively:
\begin{itemize}
    \item $\beta_\alpha$ captures how the ``belief stiffness'' (sensitivity to spatial variation) changes with scale.
    \item $\beta_\beta$ measures how strongly beliefs couple to entropy/noise at different resolutions.
    \item $\beta_\gamma$ measures how strongly beliefs amplify the gradient-aligned policy flow $\langle\nabla\Phi,\mathbf{v}\rangle$.
\end{itemize}

At one-loop (Gaussian) level, the corrections come from integrals of the form
\[
\int_{\Lambda/b < \|\mathbf{k}\| \le \Lambda}
\frac{d^d k}{(2\pi)^d}\;
\frac{1}{\omega^2 + \alpha\|\mathbf{k}\|^2}
\sim C_d\,\Lambda^{d-2}
\]
for some dimension-dependent constant $C_d$, leading to power-law or
logarithmic running depending on $d$.

\begin{proposition}[Scale-Dependent Predictive Market Regimes]
If the beta functions drive $\alpha(\ell)$, $\beta(\ell)$, or
$\gamma(\ell)$ toward fixed points as $\ell\to\infty$ (IR limit), the
futarchy dynamics exhibit scale-invariant prediction and policy
selection. Otherwise, crossovers between distinct predictive regimes
occur as one moves from micro to macro scales.
\end{proposition}

\begin{proof}
Standard RG reasoning: fixed points correspond to scale invariance of
the effective action; absence of fixed points implies flow between
regimes with different effective couplings.
\end{proof}

% ------------------------------------------------------------
\subsection*{E. Interpretation: Information Coarse-Graining in Futarchy}

The RG picture provides a natural interpretation of belief aggregation:

\begin{itemize}
    \item Integrating out $B_{>}$ corresponds to ignoring short-wavelength,
    local, or idiosyncratic speculative noise and retaining only the
    coarse-grained, robust belief structure.
    \item The scale dependence of $(\alpha,\beta,\gamma)$ describes how
    belief stiffness, noise sensitivity, and policy leverage change as we 
    move from micro-markets to macro-markets.
    \item RG fixed points correspond to \emph{stable predictive regimes}:
    scales at which the aggregate futarchy dynamics look self-similar.
\end{itemize}

This yields a renormalization-based restatement of Hanson’s core
intuition:
\begin{quote}
\emph{Prediction markets that operate across scales increasingly 
filter out local noise and converge to stable, scale-invariant
aggregated beliefs, which in turn guide macro policy.}
\end{quote}

% ------------------------------------------------------------
\subsection*{F. Toward an RG Classification of Economic Prediction Regimes}

In principle, one could classify economic environments by the structure
of their futarchy RG flows:
\begin{itemize}
    \item \textbf{Efficient-futarchy regime:} $\gamma(\ell)\to\gamma^\ast>0$
    and $\beta(\ell)\to0$ as $\ell\to\infty$, implying highly informative
    beliefs and low noise at large scales.
    \item \textbf{Noisy-futarchy regime:} $\beta(\ell)$ grows while
    $\gamma(\ell)$ flows to zero, implying that noise dominates and
    predictive coupling to policy vanishes at macro scales.
    \item \textbf{Critical-futarchy regime:} multiple couplings flow to
    a nontrivial fixed point, suggesting scale-free fluctuations in
    beliefs and policy, akin to critical phenomena.
\end{itemize}

A full analysis would derive explicit beta functions for specific
economic microfoundations and belief-update mechanisms. This appendix
establishes the conceptual backbone for such a program: futarchy, in
the RSVP field-theoretic sense, is an RG process that progressively
coarse-grains speculative information while reshaping the couplings
between belief, value, policy, and entropy fields.

% ============================================================
% Appendix: Concrete Instantiation in an Aiyagari Economy
%           (RSVP + Futarchy Specialization)
% ============================================================

\section*{Appendix: RSVP--Futarchy Formulation of an Aiyagari Economy}

This appendix specializes the RSVP--Futarchy framework to a canonical
incomplete-markets (Aiyagari) model. We explicitly define the economic
state space, map it to RSVP fields $(\Phi,\mathbf{v},S,B)$, construct
the equilibrium operator $T(\Phi,\mathbf{v},S)$, and then insert this
operator into the RSVP--Futarchy Lagrangian.

% ------------------------------------------------------------
\subsection*{A. Baseline Aiyagari Model}

Consider a continuum of ex-ante identical households indexed by $i$,
each solving:
\[
\max_{\{c_{it},a_{i,t+1}\}}
\mathbb{E}_0 \sum_{t=0}^{\infty} \beta^t u(c_{it}),
\]
subject to
\begin{align}
c_{it} + a_{i,t+1} &= (1+r_t)a_{it} + w_t z_{it},\\
a_{i,t+1} &\ge \underline{a},
\end{align}
where:
\begin{itemize}
    \item $a_{it}$: individual assets,
    \item $z_{it} \in \mathcal{Z}$: idiosyncratic productivity shock following a Markov chain,
    \item $r_t$: interest rate,
    \item $w_t$: wage rate,
    \item $u(\cdot)$: period utility (e.g.\ CRRA).
\end{itemize}

The aggregate capital stock $K_t$ is determined by the stationary
distribution of assets across households:
\[
K_t = \int a\, d\mu_t(a,z),
\]
where $\mu_t$ is the cross-sectional distribution over $(a,z)$.

Firms are competitive with production function $F(K_t,L_t)$, with
$L_t$ denoting effective labor. Prices satisfy:
\[
r_t = F_K(K_t,L_t) - \delta,\qquad
w_t = F_L(K_t,L_t).
\]

A stationary recursive competitive equilibrium consists of decision
rules, prices, and a stationary wealth distribution satisfying household
optimality, firm optimality, and market clearing.

% ------------------------------------------------------------
\subsection*{B. State Space and RSVP Fields}

We define the individual state as
\[
x = (a,z,K),
\]
where $K$ is the aggregate capital stock (or belief about it). The
economic state manifold is:
\[
\mathcal{X} = [\underline{a},\overline{a}] \times \mathcal{Z} \times [0,\overline{K}].
\]

We embed this into RSVP fields:
\begin{itemize}
    \item \textbf{Value Field $\Phi(x)$}: the household’s lifetime value
    function given $(a,z,K)$:
    \[
    \Phi(a,z,K) \approx V(a,z;K).
    \]
    \item \textbf{Policy Field $\mathbf{v}(x)$}: a vector giving the
    optimal controls. For a simple Aiyagari model,
    \[
    \mathbf{v}(a,z,K) = 
    \begin{pmatrix}
        c^\star(a,z,K) \\
        a'^\star(a,z,K)
    \end{pmatrix},
    \]
    representing consumption and next-period assets.
    \item \textbf{Entropy Field $S(x)$}: an entropy density over
    $(a,z,K)$, encoding the complexity or dispersion of the
    cross-sectional distribution and beliefs.
    \item \textbf{Belief Field $B(x)$}: a futarchy prediction field
    representing market-implied beliefs about future aggregates
    (e.g.\ $K'$, $r'$, $w'$) contingent on policies.
\end{itemize}

Thus the RSVP fields live on $\mathcal{X}$ and encode both micro and
macro structure.

% ------------------------------------------------------------
\subsection*{C. Equilibrium Operator for the Aiyagari Model}

We now construct the operator $T(\Phi,\mathbf{v},S)$ whose zeroes
characterize equilibrium.

\subsubsection*{1. Bellman Residual}

Let $\mathcal{T}$ be the Bellman operator:
\[
(\mathcal{T}\Phi)(a,z,K)
=
\max_{a'}
\left\{
    u\big((1+r(K))a + w(K)z - a'\big)
    + \beta \mathbb{E}\big[\Phi(a',z',K') \mid z\big]
\right\}.
\]

In a stationary equilibrium with constant $K$, we have $K' = K$, but in
general we allow aggregate dynamics $K' = G(\mu)$ or $K' = G(\Phi,\mathbf{v},S)$.

The value function residual is:
\[
R_V(a,z,K)
=
\Phi(a,z,K) - (\mathcal{T}\Phi)(a,z,K).
\]

\subsubsection*{2. Policy Consistency Residual}

The policy field $\mathbf{v}$ must attain the Bellman maximizer. Let:
\[
a'_\Phi(a,z,K) = 
\arg\max_{a'}
\left\{
    u\big((1+r(K))a + w(K)z - a'\big)
    + \beta \mathbb{E}[\Phi(a',z',K')\mid z]
\right\}.
\]

Define the policy residual as:
\[
R_{\mathbf{v}}(a,z,K)
=
\begin{pmatrix}
c(a,z,K) - \big((1+r(K))a + w(K)z - a'(a,z,K)\big) \\
a'(a,z,K) - a'_\Phi(a,z,K)
\end{pmatrix},
\]
where
\[
\mathbf{v}(a,z,K) =
\begin{pmatrix}
c(a,z,K) \\
a'(a,z,K)
\end{pmatrix}.
\]

\subsubsection*{3. Distribution and Market-Clearing Residual}

Let $\mu(a,z,K)$ denote a density-compatible representation of the
state distribution. The Kolmogorov (law-of-motion) residual $R_\mu$
captures the mismatch between the implied stationary distribution under
policy $\mathbf{v}$ and the actual distribution $\mu$. At a high level:
\[
R_\mu(a,z,K)
=
\mu(a,z,K)
-
\mathcal{P}_{\mathbf{v}}[\mu](a,z,K),
\]
where $\mathcal{P}_{\mathbf{v}}$ is the operator that pushes $\mu$
forward via the transition $(a,z) \to (a'(a,z,K),z')$.

The aggregate capital residual is:
\[
R_K(K,\mu)
=
K - \int a\, d\mu(a,z,K).
\]

\subsubsection*{4. Entropy Residual}

We define an entropy residual capturing the discrepancy between the
actual distribution and a low-entropy, RSVP-compatible reference:
\[
R_S(a,z,K)
=
S(a,z,K) - S_{\text{ref}}(\mu)(a,z,K),
\]
where $S_{\text{ref}}$ could be, for instance, the negative log-density
or some coarse-grained entropy functional.

\subsubsection*{5. Combined Equilibrium Operator}

Collecting residuals:
\[
T(\Phi,\mathbf{v},S)
=
\begin{pmatrix}
R_V \\
R_{\mathbf{v}} \\
R_\mu \\
R_K \\
R_S
\end{pmatrix}.
\]

An Aiyagari equilibrium in RSVP form is then:
\[
T(\Phi,\mathbf{v},S) = 0.
\]

% ------------------------------------------------------------
\subsection*{D. Specializing the RSVP--Futarchy Lagrangian}

We now insert this $T(\Phi,\mathbf{v},S)$ into the futarchy Lagrangian:
\[
\mathcal{L}
=
\mathcal{L}_{\text{dyn}}
+
\mathcal{L}_{\text{equil}}
+
\mathcal{L}_{\text{belief}}
+
\mathcal{L}_{\text{futarchy}}
+
\mathcal{L}_{\text{entropy}}.
\]

\subsubsection*{1. Equilibrium Term}

The equilibrium constraint becomes:
\[
\mathcal{L}_{\text{equil}}
=
- \frac{\lambda}{2}
\left(
    |R_V|^2
  + \|R_{\mathbf{v}}\|^2
  + |R_\mu|^2
  + |R_K|^2
  + |R_S|^2
\right),
\]
integrated over $(a,z,K)$ and time.

\subsubsection*{2. Futarchy \& Belief Terms in Aiyagari Context}

In this model, the belief field $B$ can be interpreted as a prediction
market on future aggregate quantities (e.g.\ $K'$, $r'$, $w'$) and
associated welfare. For concreteness, let $B$ encode beliefs about the
expected value of some welfare functional $W[\Phi,\mu]$, such as
aggregate utility or consumption.

The belief Lagrangian remains:
\[
\mathcal{L}_{\text{belief}}
=
- \frac{\alpha}{2}\|\nabla B\|^2
- \beta B S,
\]
interpreting $\nabla$ as gradients in $(a,z,K)$.

The futarchy coupling term now aligns policy with value gradients in the
Aiyagari state space:
\[
\mathcal{L}_{\text{futarchy}}
=
\gamma\, B\, \langle \nabla \Phi,\ \tilde{\mathbf{v}} \rangle
- \frac{\kappa}{2}\|\mathbf{v} - \mathbf{v}_B\|^2.
\]

Here:
\begin{itemize}
    \item $\nabla \Phi$ is the gradient of the value function in $(a,z,K)$,
    \item $\tilde{\mathbf{v}}$ is an embedding of the policy vector into
    the tangent space of $\mathcal{X}$,
    \item $\mathbf{v}_B$ is the policy implied by the belief field $B$
    (e.g.\ the policy that maximizes predicted aggregate welfare).
\end{itemize}

Economically, this says:
\begin{quote}
Policies are driven to align with the value gradient in state space, as
weighted by market beliefs, and penalized for deviating from the
belief-implied optimal policies.
\end{quote}

\subsubsection*{3. Entropy Term for the Wealth Distribution}

The entropy term becomes:
\[
\mathcal{L}_{\text{entropy}}
=
- \eta |\nabla S|^2
- \xi S \log S,
\]
where $S$ is related to the cross-sectional density $\mu$ (e.g.\
$S = -\log \mu$ or a smoothed version thereof). This penalizes highly
fragmented or unstable distributions and favors smooth, coherent
wealth distributions consistent with RSVP lamphrodynamics.

% ------------------------------------------------------------
\subsection*{E. Interpretation in the Aiyagari Context}

In this concrete setting, the RSVP--Futarchy theory tells the following
story:
\begin{itemize}
    \item $\Phi(a,z,K)$: the household value function and, more broadly,
    a welfare potential over the microstate space.
    \item $\mathbf{v}(a,z,K)$: policy rules mapping from states to
    controls $(c,a')$.
    \item $S(a,z,K)$: an entropy field encoding the dispersion and
    complexity of the wealth and shock distribution.
    \item $B(a,z,K)$: a predictive field encoding market beliefs about
    how different states and policies affect future aggregates and
    welfare.
\end{itemize}

The equilibrium operator $T(\Phi,\mathbf{v},S)$ enforces:
\begin{enumerate}
    \item Bellman consistency,
    \item policy optimality,
    \item distributional stationarity,
    \item market clearing,
    \item entropy coherence.
\end{enumerate}

The futarchy coupling terms ensure that:
\begin{enumerate}
    \item beliefs $B$ aggregate dispersed information about how
    $(\Phi,\mathbf{v},S)$ shape future states;
    \item policy fields $\mathbf{v}$ are driven toward those that
    maximize predicted welfare under $B$;
    \item the resulting dynamics form a lamphrodynamic flow on the
    space of Aiyagari equilibria, selecting stable, low-entropy,
    belief-consistent configurations.
\end{enumerate}

Thus a classical incomplete-markets Aiyagari model is realized as a
particular instance of the RSVP--Futarchy field theory, with the
standard equilibrium conditions emerging as the zero-gradient limit of a
more general lamphrodynamic and variational structure.

% ============================================================
% Appendix: Stationary RSVP Free-Energy for Aiyagari Equilibrium
% ============================================================

\section*{Appendix: Stationary RSVP Free-Energy\\
for an Aiyagari Equilibrium}

In this appendix we specialize the RSVP--Futarchy formulation of the
Aiyagari economy to a stationary equilibrium and construct an explicit
free-energy functional over the state space. The goal is to make
concrete the notion that a stationary Aiyagari equilibrium corresponds
to a minimizer of a lamphrodynamic free energy.

% ------------------------------------------------------------
\subsection*{A. Stationary Setup}

We restrict attention to time-invariant configurations:
\[
\partial_t \Phi = 0,\quad
\partial_t \mathbf{v} = 0,\quad
\partial_t S = 0,\quad
\partial_t B = 0.
\]

The state space is
\[
x = (a,z,K),\qquad
\mathcal{X} = [\underline{a},\overline{a}] \times \mathcal{Z} \times [0,\overline{K}],
\]
and we assume that the aggregate capital stock is constant at
equilibrium, $K=K^\ast$, so that effectively the state is $(a,z)$ and
$K^\ast$ enters as a parameter.

We write the fields as:
\[
\Phi(a,z),\quad
\mathbf{v}(a,z) = \begin{pmatrix} c(a,z) \\ a'(a,z) \end{pmatrix},\quad
S(a,z),\quad
B(a,z),
\]
with $K^\ast$ implicitly determined by the stationary wealth
distribution.

% ------------------------------------------------------------
\subsection*{B. Stationary Equilibrium Operator}

The equilibrium operator $T(\Phi,\mathbf{v},S)$ defined previously
simplifies under stationarity and constant $K^\ast$.

\subsubsection*{1. Bellman Residual}

The Bellman operator becomes:
\[
(\mathcal{T}\Phi)(a,z)
=
\max_{a'}
\left\{
    u\big((1+r^\ast)a + w^\ast z - a'\big)
    + \beta \mathbb{E}\big[\Phi(a',z')\mid z\big]
\right\},
\]
where $r^\ast = r(K^\ast)$ and $w^\ast = w(K^\ast)$.

The value residual is:
\[
R_V(a,z)
=
\Phi(a,z) - (\mathcal{T}\Phi)(a,z).
\]

\subsubsection*{2. Policy Residual}

Define
\[
a'_\Phi(a,z)
=
\arg\max_{a'}
\left\{
    u\big((1+r^\ast)a + w^\ast z - a'\big)
    + \beta \mathbb{E}\big[\Phi(a',z')\mid z\big]
\right\}.
\]

The policy residual remains:
\[
R_{\mathbf{v}}(a,z)
=
\begin{pmatrix}
c(a,z) - \big((1+r^\ast)a + w^\ast z - a'(a,z)\big) \\
a'(a,z) - a'_\Phi(a,z)
\end{pmatrix}.
\]

\subsubsection*{3. Distribution and Market-Clearing Residuals}

Let $\mu(a,z)$ denote the stationary distribution. The law-of-motion
residual becomes:
\[
R_\mu(a,z)
=
\mu(a,z) - \mathcal{P}_{\mathbf{v}}[\mu](a,z),
\]
where $\mathcal{P}_{\mathbf{v}}[\mu]$ is the stationary fixed point of
the Markov operator induced by $(a'(a,z),z')$.

The capital residual simplifies to:
\[
R_K
=
K^\ast - \int a\, d\mu(a,z).
\]

\subsubsection*{4. Entropy Residual}

We define a reference entropy $S_{\text{ref}}(\mu)$, for example
\[
S_{\text{ref}}(a,z) = -\log \mu(a,z)
\]
(up to normalization). The entropy residual is:
\[
R_S(a,z)
=
S(a,z) - S_{\text{ref}}(a,z).
\]

\subsubsection*{5. Combined Stationary Operator}

Thus the stationary equilibrium operator is:
\[
T(\Phi,\mathbf{v},S)
=
\begin{pmatrix}
R_V(a,z) \\
R_{\mathbf{v}}(a,z) \\
R_\mu(a,z) \\
R_K \\
R_S(a,z)
\end{pmatrix},
\]
and a stationary equilibrium is a zero of $T$.

% ------------------------------------------------------------
\subsection*{C. Stationary Free-Energy Functional}

We now define the stationary RSVP free-energy functional:
\begin{equation}
\mathcal{E}[\Phi,\mathbf{v},S;\mu,K^\ast]
=
\frac{1}{2}
\int_{\underline{a}}^{\overline{a}}
\sum_{z\in\mathcal{Z}}
w(a,z)
\Big(
    |R_V(a,z)|^2
  + \|R_{\mathbf{v}}(a,z)\|^2
  + |R_\mu(a,z)|^2
  + |R_S(a,z)|^2
\Big)\, da
+ \frac{\omega}{2} |R_K|^2,
\end{equation}
where $w(a,z)$ is a weighting function (e.g.\ the stationary density
$\mu(a,z)$) and $\omega>0$ is a weight on the capital residual.

\begin{theorem}[Stationary Aiyagari Equilibrium as Free-Energy Minimizer]
Suppose the functional $\mathcal{E}$ is differentiable and the
equilibrium operator $T$ is nondegenerate in the sense described in the
main text. Then:
\[
(\Phi^\ast,\mathbf{v}^\ast,S^\ast,\mu^\ast,K^\ast)
\text{ is a stationary Aiyagari equilibrium}
\iff
(\Phi^\ast,\mathbf{v}^\ast,S^\ast,\mu^\ast,K^\ast)
\text{ minimizes } \mathcal{E}.
\]
\end{theorem}

\begin{proof}
If $T(\Phi^\ast,\mathbf{v}^\ast,S^\ast)=0$ and $R_K=0$, then each
residual vanishes, implying $\mathcal{E}=0$, the global minimum. 
Conversely, if $(\Phi^\ast,\mathbf{v}^\ast,S^\ast,\mu^\ast,K^\ast)$ 
minimizes $\mathcal{E}$, the first variation must vanish. Under the 
nondegeneracy assumption, this implies each residual is zero, i.e. 
$T=0$ and $R_K=0$.
\end{proof}

This makes precise the statement that the stationary equilibrium of an
Aiyagari economy is a lamphrodynamically relaxed configuration in RSVP
field space.

% ------------------------------------------------------------
\subsection*{D. Stationary Lamphrodynamic Gradient Flow}

Even in the stationary setting, one can define an abstract gradient flow
in a fictitious ``relaxation time'' $\tau$:
\begin{align}
\frac{\partial \Phi}{\partial \tau}
&= - \frac{\delta \mathcal{E}}{\delta \Phi},\\[4pt]
\frac{\partial \mathbf{v}}{\partial \tau}
&= - \frac{\delta \mathcal{E}}{\delta \mathbf{v}},\\[4pt]
\frac{\partial S}{\partial \tau}
&= - \frac{\delta \mathcal{E}}{\delta S},\\[4pt]
\frac{\partial \mu}{\partial \tau}
&= - \frac{\delta \mathcal{E}}{\delta \mu},\\[4pt]
\frac{\partial K^\ast}{\partial \tau}
&= - \frac{\partial \mathcal{E}}{\partial K^\ast}.
\end{align}

This flow monotonically decreases $\mathcal{E}$ and converges, under
suitable conditions, to the stationary equilibrium. In numerical deep
learning solvers, discrete steps of stochastic gradient descent on
network parameters approximate this continuum lamphrodynamic flow.

% ============================================================
% Appendix: RSVP--Futarchy Formulation of a HANK New Keynesian Model
% ============================================================

\section*{Appendix: RSVP--Futarchy Formulation\\
of a HANK-Style New Keynesian Economy}

In this appendix we sketch a Heterogeneous-Agent New Keynesian (HANK)
economy in RSVP--Futarchy terms. The goal is to show how macro policy
rules (monetary and/or fiscal), shocks, and heterogeneous households fit
naturally into the $(\Phi,\mathbf{v},S,B)$ field-theoretic framework.

% ------------------------------------------------------------
\subsection*{A. Stylized HANK New Keynesian Setup}

We consider a stylized HANK environment with:
\begin{itemize}
    \item heterogeneous households indexed by assets $a$, income states $z$;
    \item aggregate variables: output gap $x_t$, inflation $\pi_t$,
          nominal interest rate $i_t$;
    \item a monetary policy rule and possibly fiscal instruments;
    \item idiosyncratic and aggregate shocks.
\end{itemize}

For concreteness, suppose aggregate dynamics obey (in log-linear form):
\begin{align}
x_t &= \mathbb{E}_t x_{t+1} - \sigma (i_t - \mathbb{E}_t \pi_{t+1} - r_t^{n}),\\
\pi_t &= \beta \mathbb{E}_t \pi_{t+1} + \kappa x_t + u_t,
\end{align}
where $u_t$ is a cost-push shock and $r_t^{n}$ the natural real rate.

A Taylor-type monetary rule:
\[
i_t = \bar{i} + \phi_\pi \pi_t + \phi_x x_t + \varepsilon_t,
\]
with $\varepsilon_t$ a monetary policy shock.

Heterogeneous households face incomplete markets and solve problems
similar in spirit to the Aiyagari households, but in a time-varying
aggregate environment.

% ------------------------------------------------------------
\subsection*{B. State Space and RSVP Fields in HANK--NK}

We define a combined micro--macro state:
\[
x = (a,z,x^A),
\qquad
x^A = (x,\pi,i,\dots),
\]
where $x^A$ collects aggregate variables (output gap, inflation, rate,
and possibly others).

The state manifold is:
\[
\mathcal{X}
=
[a_{\min},a_{\max}]
\times \mathcal{Z}
\times \mathcal{X}^A,
\]
where $\mathcal{X}^A$ is the macro-state space.

We define RSVP fields as:
\begin{itemize}
    \item \textbf{Value Field $\Phi(a,z,x^A)$}: HANK value function,
    incorporating both micro and macro states.
    \item \textbf{Policy Field $\mathbf{v}(a,z,x^A)$}: a vector of
    household controls and possibly macro policy levers. For instance:
    \[
    \mathbf{v} =
    \begin{pmatrix}
        c(a,z,x^A) \\
        a'(a,z,x^A)
    \end{pmatrix}
    \quad
    \text{(micro)},
    \]
    plus
    \[
    \mathbf{v}^{\text{pol}}(x^A)
    =
    \begin{pmatrix}
        i(x^A) \\
        f(x^A)
    \end{pmatrix}
    \quad
    \text{(macro policy rule: monetary $i$, fiscal $f$)}.
    \]
    For brevity, we write the total policy field as $\mathbf{v}$,
    with micro and macro components.
    \item \textbf{Entropy Field $S(a,z,x^A)$}: encoding the complexity of
    the joint distribution over assets, idiosyncratic states, and macro
    configurations.
    \item \textbf{Belief Field $B(a,z,x^A)$}: a futarchy belief field
    representing prediction-market beliefs over future aggregates and
    welfare outcomes given candidate policy rules.
\end{itemize}

% ------------------------------------------------------------
\subsection*{C. Equilibrium Operator in HANK--NK}

We sketch the main components of $T(\Phi,\mathbf{v},S)$.

\subsubsection*{1. Micro Bellman Residual}

Households solve:
\[
\Phi(a,z,x^A)
=
\max_{a'}
\left\{
    u(c(a,z,x^A))
    + \beta \mathbb{E}
    \big[
        \Phi(a',z',x^{A\prime})
    \mid a,z,x^A
    \big]
\right\},
\]
where $c$ obeys budget constraints given prices determined by
$x^A=(x,\pi,i,\dots)$ and policy $\mathbf{v}^{\text{pol}}$.

We define:
\[
R_V(a,z,x^A)
=
\Phi(a,z,x^A) - (\mathcal{T}\Phi)(a,z,x^A),
\]
where $\mathcal{T}$ is the corresponding Bellman operator.

\subsubsection*{2. Micro Policy Residual}

As in the Aiyagari case, define an optimal $a'_\Phi$ and construct:
\[
R_{\mathbf{v}}^{\text{micro}}(a,z,x^A)
=
\begin{pmatrix}
c(a,z,x^A) - \text{budget-consistent}(a,z,x^A,a') \\
a'(a,z,x^A) - a'_\Phi(a,z,x^A)
\end{pmatrix}.
\]

\subsubsection*{3. Macro NK Residuals}

We define macro residuals for the NK block:
\begin{align}
R_x(x^A)
&=
x - \big[
    \mathbb{E}x' - \sigma (i - \mathbb{E}\pi' - r^n)
\big],\\
R_\pi(x^A)
&=
\pi - \big[
    \beta \mathbb{E}\pi' + \kappa x + u
\big],\\
R_i(x^A)
&=
i - \big[\bar{i} + \phi_\pi\pi + \phi_x x + \varepsilon\big].
\end{align}

These can be collected in a macro residual vector:
\[
R_{\text{NK}}(x^A)
=
\begin{pmatrix}
R_x(x^A) \\
R_\pi(x^A) \\
R_i(x^A)
\end{pmatrix}.
\]

\subsubsection*{4. Distribution and Entropy Residuals}

Let $\mu(a,z,x^A)$ be the joint distribution over micro and macro
states (or a conditional distribution over $(a,z)$ given $x^A$ and an
aggregate law of motion for $x^A$). We define:
\[
R_\mu(a,z,x^A)
=
\mu(a,z,x^A) - \mathcal{P}_{\mathbf{v}}[\mu](a,z,x^A),
\]
and a reference entropy $S_{\text{ref}}(\mu)$, with:
\[
R_S(a,z,x^A)
=
S(a,z,x^A) - S_{\text{ref}}(a,z,x^A).
\]

\subsubsection*{5. Combined HANK--NK Operator}

Collecting all components:
\[
T(\Phi,\mathbf{v},S)
=
\begin{pmatrix}
R_V(a,z,x^A) \\
R_{\mathbf{v}}^{\text{micro}}(a,z,x^A) \\
R_{\text{NK}}(x^A) \\
R_\mu(a,z,x^A) \\
R_S(a,z,x^A)
\end{pmatrix}.
\]

A HANK--NK equilibrium satisfies $T(\Phi,\mathbf{v},S)=0$.

% ------------------------------------------------------------
\subsection*{D. Futarchy Coupling for Macro Policy Rules}

In a futarchy setting, prediction markets bet on the consequences of
alternative macro policy rules. Let $\theta$ denote rule coefficients
(e.g.\ $\phi_\pi,\phi_x$, or parameters of a richer monetary/fiscal
rule), and encode them as part of the macro policy field
$\mathbf{v}^{\text{pol}}$.

The futarchy term in the Lagrangian is adapted to:
\[
\mathcal{L}_{\text{futarchy}}^{\text{HANK--NK}}
=
\gamma\, B(a,z,x^A)\,
    \langle \nabla \Phi(a,z,x^A),\ \tilde{\mathbf{v}}(a,z,x^A) \rangle
- \frac{\kappa}{2}
    \|\mathbf{v}^{\text{pol}}(x^A) - \mathbf{v}^{\text{pol}}_B(x^A)\|^2.
\]

Here:
\begin{itemize}
    \item $B$ aggregates predictions over future welfare given policy
    rule parameters $\theta$,
    \item $\mathbf{v}^{\text{pol}}_B$ is the policy rule that maximizes
    expected welfare under beliefs $B$,
    \item $\tilde{\mathbf{v}}$ is the embedding of micro/macro policy
    directions into the tangent space of $\mathcal{X}$.
\end{itemize}

This implements Hanson’s principle:
\[
\text{choose policy rules that maximize predicted welfare}.
\]

% ------------------------------------------------------------
\subsection*{E. HANK--NK Free-Energy Functional}

We can now define a HANK--NK RSVP free-energy functional:
\begin{align}
\mathcal{E}_{\text{HANK--NK}}[\Phi,\mathbf{v},S]
&=
\frac{1}{2}
\int_{\mathcal{X}}
\Big(
    |R_V|^2
  + \|R_{\mathbf{v}}^{\text{micro}}\|^2
  + |R_\mu|^2
  + |R_S|^2
\Big)\,d\mu_x
\\[4pt]
&\quad
+ \frac{1}{2}
\int_{\mathcal{X}^A}
\|R_{\text{NK}}(x^A)\|^2\, d\mu_A,
\end{align}
where $d\mu_x$ and $d\mu_A$ are appropriate measures over the full
state space and macro space, respectively.

\begin{theorem}[HANK--NK Equilibrium as RSVP Free-Energy Minimizer]
Under regularity and nondegeneracy assumptions analogous to those in
the stationary Aiyagari case,
\[
T(\Phi^\ast,\mathbf{v}^\ast,S^\ast)=0
\quad\Longleftrightarrow\quad
(\Phi^\ast,\mathbf{v}^\ast,S^\ast)
\text{ minimizes }
\mathcal{E}_{\text{HANK--NK}}.
\]
\end{theorem}

\begin{proof}
Identical in structure to the Aiyagari free-energy theorem: if all
residuals vanish, $\mathcal{E}_{\text{HANK--NK}}=0$, the global minimum.
Conversely, a minimizer must have vanishing first variation, which
implies each residual is zero under nondegeneracy.
\end{proof}

Thus a HANK--NK equilibrium emerges as a lamphrodynamically relaxed
configuration of the RSVP fields, with futarchy couplings directing the
macro policy rule toward belief-weighted welfare optima.

% ------------------------------------------------------------
\subsection*{F. Interpretation}

In summary:
\begin{itemize}
    \item The \emph{micro} side (heterogeneous households) is captured by
    $\Phi(a,z,x^A)$, $\mathbf{v}^{\text{micro}}$, and $S(a,z,x^A)$.
    \item The \emph{macro} side (NK block and policy rules) is captured
    by $\mathbf{v}^{\text{pol}}(x^A)$, $R_{\text{NK}}(x^A)$, and their
    contribution to the free energy.
    \item The \emph{belief} field $B(a,z,x^A)$ integrates predictive
    information about the macro effects of policy rules, and, once
    integrated out, induces nonlocal couplings in the effective action.
    \item Equilibrium is the joint minimum of the free-energy functional
    over both micro and macro fields, respecting futarchy’s separation
    of values (encoded in $\Phi$) and beliefs (encoded in $B$).
\end{itemize}

This closes the loop between HANK microfoundations, New Keynesian
aggregate dynamics, deep learning–style lamphrodynamic solvers, and
Hanson-style futarchy within a single RSVP field-theoretic framework.


\end{document}
